%====================================
% KAPITEL 4: ERGEBNISSE DER LITERATURANALYSE
%====================================
\section{Ergebnisse der Literaturanalyse}

Das folgende Kapitel pr\"asentiert die Ergebnisse der systematischen Literaturanalyse. Zun\"achst wird das finale Sample deskriptiv ausgewertet, um die Struktur der Befundlage -- Publikationsverlauf, Untersuchungsmethoden und Untersuchungskontexte -- transparent zu machen und die Belastbarkeit der nachfolgenden inhaltlichen Aussagen einordnen zu k\"onnen (Kapitel~4.1). Anschlie\ss{}end werden die identifizierten Change-Management-Faktoren kategorisiert und entlang des Kategoriensystems aus Kapitel~3.3 dargestellt (Kapitel~4.2). Kapitel~4.3 verkn\"upft die Faktoren mit den Erfolgsdimensionen des DeLone-\&-McLean-Modells und \"uberpr\"uft damit den vorl\"aufigen Bezugsrahmen aus Kapitel~2.4.

\subsection{Deskriptive Auswertung des Samples}

\textit{[GER\"UST -- Zahlen und Aussagen nach durchgef\"uhrter Suche einsetzen. Alle Platzhalter sind mit $\square$ markiert.]}

\textbf{Zusammensetzung des Samples}

Die Datenbanksuche ergab insgesamt $\square$ Treffer, von denen nach Entfernung von $\square$ Dubletten $\square$ Beitr\"age in das Titel-Abstract-Screening eingingen. Nach Anwendung der Ein- und Ausschlusskriterien (Kapitel~3.3) verblieben $\square$ Beitr\"age f\"ur die Volltextpr\"ufung, von denen $\square$ in das finale Sample eingeschlossen wurden. Die R\"uckw\"arts- und Vorw\"artssuche erg\"anzte das Sample um $\square$ weitere Beitr\"age, sodass der Analyse insgesamt $\square$ Beitr\"age zugrunde liegen. Der vollst\"andige Auswahlprozess ist in Abbildung~2 (Kapitel~3.3) dokumentiert; die h\"aufigsten Ausschlussgr\"unde auf Volltextebene waren $\square$ \textit{[z.\,B.: fehlender Change-Management-Bezug, rein technischer Fokus, kein abgrenzbarer KMU-Kontext]}.

\textbf{Publikationsverlauf und Publikationsorgane}

Betrachtet man die zeitliche Verteilung der Beitr\"age (Abbildung~3), zeigt sich $\square$ \textit{[z.\,B.: ein weitgehend konstantes Publikationsaufkommen / ein deutlicher Anstieg ab dem Jahr $\square$, der mit der zunehmenden Verbreitung von Cloud-ERP-L\"osungen im Mittelstand zusammenfallen d\"urfte]}. \textit{[Abbildung 3: Balkendiagramm Beitr\"age je Erscheinungsjahr 2015--$\square$; Eigene Darstellung]}

Die Beitr\"age verteilen sich auf $\square$ verschiedene Publikationsorgane. Am h\"aufigsten vertreten sind $\square$ \textit{[Zeitschriften/Konferenzen nennen]}. $\square$ \textit{[Einordnender Satz, z.\,B.: Die Streuung \"uber Zeitschriften der Informationssystemforschung, des Operations Management und der KMU-Forschung best\"atigt den in Kapitel~1.1 konstatierten fragmentierten Charakter des Forschungsfelds.]}

\textbf{Untersuchungsmethoden}

Methodisch dominieren $\square$ \textit{[z.\,B.: Fallstudien (n = $\square$), gefolgt von quantitativen Befragungen (n = $\square$) und Literaturanalysen (n = $\square$)]} (Tabelle~3). $\square$ \textit{[Einordnung, z.\,B.: Der hohe Anteil an Einzel- und Mehrfachfallstudien ist f\"ur die KMU-Forschung typisch, begrenzt jedoch die statistische Generalisierbarkeit der Befunde -- ein Punkt, der in Kapitel~6.2 aufgegriffen wird.]}

\begin{table}[H]
\centering
\begin{tabularx}{0.6\textwidth}{lcc}
\toprule
\textbf{Methode} & \textbf{Anzahl} & \textbf{Anteil in \%} \\
\midrule
$\square$ & $\square$ & $\square$ \\
$\square$ & $\square$ & $\square$ \\
$\square$ & $\square$ & $\square$ \\
\bottomrule
\end{tabularx}
\caption{Sample nach Untersuchungsmethode (Eigene Darstellung)}
\label{tab:untersuchungsmethoden}
\end{table}

\textbf{Untersuchungskontexte}

Hinsichtlich der geografischen Verteilung stammen die Beitr\"age aus $\square$ \textit{[Regionen/L\"ander mit n nennen]}. $\square$ \textit{[Falls zutreffend: Auff\"allig ist der geringe Anteil an Studien aus dem deutschsprachigen Raum (n = $\square$), obwohl der Mittelstandsbegriff hier besonders gepr\"agt ist -- die \"Ubertragbarkeit internationaler KMU-Befunde auf den deutschen Mittelstand wird in Kapitel~5.2 reflektiert.]}

Bezogen auf die Unternehmensgr\"o\ss{}e untersuchen $\square$ Beitr\"age explizit KMU beziehungsweise mittelst\"andische Unternehmen, w\"ahrend $\square$ Beitr\"age gr\"o\ss{}en\"ubergreifende Samples mit abgrenzbarer KMU-Teilstichprobe aufweisen und $\square$ Beitr\"age aufgrund \"ubertragbarer Befunde ohne expliziten Gr\"o\ss{}enfokus eingeschlossen wurden \textit{[letzteres nur, falls das weiche Kontextkriterium aus Tabelle~2 genutzt wurde]}. $\square$ \textit{[Ggf. Branchenverteilung erg\"anzen, z.\,B. Dominanz des produzierenden Gewerbes.]}

\textbf{Zwischenfazit zur Befundlage}
